\def\stat{korepanov}





\def\tit{НЕКОТОРЫЕ ПОДХОДЫ К~РАЗРАБОТКЕ ТЕХНОЛОГИЙ ТОНКОГО КЛИЕНТА 


ДЛЯ~ЗАЩИЩЕННЫХ ИНФОРМАЦИОННЫХ СИСТЕМ}





\def\titkol{Некоторые подходы к~разработке технологий тонкого клиента} 


%для~защищенных информационных систем} 





\def\autkol{Э.\,Р.~Корепанов}





\def\aut{Э.\,Р.~Корепанов$^1$}





\titel{\tit}{\aut}{\autkol}{\titkol}





%{\renewcommand{\thefootnote}{\fnsymbol{footnote}}\footnotetext[1]


%{Работа выполнена при поддержке РФФИ (грант 11-07-00112).}}





\renewcommand{\thefootnote}{\arabic{footnote}}


\footnotetext[1]{Институт проблем информатики Российской академии наук, ekorepanov@ipiran.ru}





\vspace*{6pt}


 





\Abst{Рассматриваются вопросы развития технологий тонкого клиента для российских 


защищенных информационных систем, основанные на мировом опыте создания 


инфраструктуры виртуализации персональных компьютеров (VDI) и использовании 


аппаратных супертонких клиентов.}





\vspace*{2pt}





\KW{защищенные информационные системы; технология тонкого клиента; 


инфраструктура виртуализации персональных компьютеров; персональный компьютер по 


протоколу IP}





\vspace*{6pt}





\vskip 14pt plus 9pt minus 6pt








      \thispagestyle{headings}





      


            \label{st\stat}


            


            \vspace*{-12pt}





\section{Введение}





       За последние 30~лет технологии тонкого клиента в России прошли путь от 


доминирующего средства работы конечного пользователя с большими ЭВМ в середине 


1980-х гг.\ до почти полного забвения в конце 1980-х и начале 1990-х~гг.--- 


эпоху расцвета персональных компьютеров~--- и снова успешно возродились в новом 


тысячелетии благодаря интра\-нет-тех\-но\-ло\-ги\-ям. В~зависимости от вычислительной 


мощности серверной инфраструктуры, пропускной способности каналов связи, 


графических возможностей абонентских средств использовались те или иные варианты 


технологии тонкого клиента, наиболее подходящие для решения конкретных прикладных 


задач. В~целом централизованная модель вычислений, реализованная с применением 


технологий тонкого клиента, оставалась достаточно постоянной на протяжении 


десятилетия, до начала 2000-х~гг., и подробно представлена в~[1].


       


       Когда речь идет о корпоративных информационных системах, можно при\-вес\-ти не 


менее девяти основных преимуществ использования технологий тонкого клиента.


       \begin{enumerate}[1.]


\item       \textit{Надежность}~--- использование аппаратуры сервера виртуализации, 


рассчитанной на круглосуточное непрерывное функционирование, повышает надежность 


работы и хранения данных. Выход из строя терминала пользователя не влечет за собой 


потерю данных на сервере.\\[-10pt]


     \item  


       \textit{Централизация}~--- все данные хранятся только на серверах, что упрощает 


для администраторов процедуру резервного копирования, контроль версий программного 


обеспечения (ПО) и контроль доступа пользователей.\\[-10pt]


      \item 


       \textit{Безопасность данных}~--- отсутствие на пользовательских терминалах 


жестких дисков и приводов оптических дисков позволяет существенно снизить 


возможности несанкционированного копирования и выноса данных. Отсутствие 


непосредственной передачи данных (передается информация об изображении для 


монитора) по сети исключает их комплексный перехват. При краже или изъятии 


терминала нет риска потерять важные данные. Существенно снижается риск заражения 


системы вирусами.\\[-10pt]


       \item


       \textit{Эффективность}~--- загрузка многоядерного процессора на современном 


офисном компьютере в большинстве случаев не превышает 5\%. Терминальная система с 


максимальной пользой задействует вычислительные ресурсы сервера виртуализации, 


распределяя их между работающими в данный момент пользователями.\\[-10pt]


       \item


       \textit{Бесшумная работа и снижение энергопотребления на рабочем месте 


пользователя}~--- терминальные клиенты не требуют больших вычислительных ресурсов, 


поэтому используют экономичные процессоры и чипсеты системной логики с низким 


тепловыделением, что не требует мощных вентиляторов и обеспечивает бесшумную 


работу. Также обычно не требуется индивидуального источника бесперебойного питания.\\[-10pt]


       \item


       \textit{Быстрое развертывание и обновление приложений}~--- при использовании 


терминальных клиентов новое ПО устанавливается администратором только на серверах 


и сразу же становится доступным пользователям.\\[-10pt] 


       \item


       \textit{Упрощение поддержки конечных узлов}~--- при использовании 


терминальных клиентов администратору нет необходимости обходить пользовательские 


устройства для установки и настройки программ. При выходе терминального клиента из 


строя его легко заменить.\\[-10pt] 


       \item


       \textit{Работа пользователя в любом месте организации}~--- пользователь может 


подключиться к своему терминальному серверу в любом месте, где установлен 


терминальный клиент.\\[-10pt]


       \item


       \textit{Сокращение совокупной стоимости владения (Total cost of ownership, 


ТСО)}~--- уменьшение стоимости эксплуатации рабочего места за счет централизованного 


администрирования, сокращения затрат рабочего времени на обслуживание рабочих мест 


пользователей, уменьшения затрат на обучение сотрудников, сокращения 


энергопотребления, отсутствия необходимости в источниках бесперебойного питания для 


терминалов (отключение терминала\linebreak\vspace*{-12pt}





\pagebreak





\noindent


 не влияет на сервер), сокращения затрат на 


модернизацию (модернизируется только сервер) и~т.\,п. 


\end{enumerate}


       


       Рассмотрим некоторые известные варианты технологий тонкого клиента для 


защищенных информационных систем.


       


       Фирма <<Анкад>> предлагает программно-аппаратный комплекс <<Криптон-


ТК>>, предназначенный для предотвращения несанкционированного доступа к 


терминалам и серверам стандартной инфраструктуры <<тонкий клиент>>, работающим 


под управлением операционной системы Microsoft Windows Server 2003 и терминального 


сервиса Citrix MetaFrame Presentation Server~4.0. Защищенный терминал комплекса 


<<Крип\-тон-ТК>> базируется на стандартном персональном компьютере c процессором 


Intel, оснащенном специализированными ап\-па\-рат\-но-про\-грам\-мны\-ми средствами защиты, в 


том числе аппаратным модулем доверенной загрузки <<Крип\-тон-За\-мок>> и сетевой 


криптоплатой линейки <<Криптон AncNet>>. В~начале каждого сеанса работы после 


идентификации и аутентификации пользователя на такой терминал осуществляется 


доверенная загрузка с сервера необходимого для работы ПО. Эта 


технология позволяет создать дополнительный барьер защиты от несанкционированного 


доступа, не базируясь на стандартных средствах защиты операционных систем.


       


       Компания <<СВЕМЕЛ>> почти 10~лет предлагает комплекс защищенного 


терминального доступа, созданный на базе защищенных технологий корпорации Sun~--- 


SunRay-сер\-ве\-ров и терминалов, а также доверенной операционной системы 


       <<Цир\-кон~10>> (доработанной ОС Solaris~10 Update~4 с установленным Solaris 


Trusted Extensions).


       


       Программные решения для организации работы защищенного тонкого клиента на 


основе веб-тех\-но\-ло\-гий предлагаются <<НПО РусБИТех>> в виде операционной 


системы специального назначения Astra Linux Special Edition.


       


       К сожалению, указанные варианты реализации защищенных технологий тонкого 


клиента базируются на решениях, популярных в начале и середине прошлого десятилетия. 


С~тех пор получили широкое распространение многоядерные и многопроцессорные 


средства централизованной обработки данных, появились очень эффективные технологии 


виртуализации, стали доступными высокоскоростные каналы связи. Новейшие 


технологии тонкого клиента консолидировались в новое популярное направление 


       IT-ин\-дуст\-рии~--- <<облачные вычисления>> (cloud computing). 


       


       Анализ ключевых элементов, предлагаемых в коммерческом секторе 


современных технологий, позволяет выработать рациональные подходы к внедрению в 


защищенных информационных системах перспективных решений на основе 


инфраструктуры виртуализации персональных компьютеров (Virtual Desktop Infrastructure, 


VDI), аппаратных супертонких клиентов (hardware zero client) и специальных протоколов 


взаимодействия с пользователем.





%\vspace*{-24pt}





\section{VDI как основа современных технологий тонкого клиента}


       


       VDI подразумевает консолидацию информационных и вычислительных ресурсов 


пользователей в группе серверов виртуализации. При этом конечный терминал 


пользователя служит лишь для ввода управляющих сигналов с клавиатуры и мыши и 


вывода графических данных, передаваемых на него по специальным протоколам с сервера 


виртуализации.


       


       В настоящее время существует мощнейшая конкуренция программных платформ 


VDI между крупными традиционными производителями терминальных решений и 


новыми <<игроками>> на этом рынке. Среди популярных в нашей стране продуктов 


можно отметить:


       \begin{itemize}


       \item VMWare View 5;


       \item Citrix XenDesktop;


       \item Microsoft VDI на базе Hyper-V, операционной системы Windows Server (2008 


R2, 2012) и System Center Virtual Machine Manager;


       \item Red Hat Enterprise Virtualization for Desktops.


       \end{itemize}


       


       Менее распространены решения Oracle VDI, в том числе на базе унаследованных 


от Sun платформ~--- Solaris Operating System, Sun Ray Software, Sun Ray Thin Client.


       


       Необходимо отметить, что из-за кажущейся простоты технологии тонкого 


клиента существует устойчивое представление о низкой стоимости аппаратных средств 


конечного пользователя~--- терминала тонкого клиента. Однако это не совсем верно, так 


как б$\acute{\mbox{о}}$льшая часть тонких клиентов представляет собой компактный персональный 


компьютер, оснащенный полноценным центральным процессором, например Intel Atom, 


оперативной памятью от 1 до 2~ГБ, графическим чипсетом и полным набором 


интерфейсов (USB, DVI-D, Audio, Ethernet 1~Гбит/с, PS/2). Наиболее существенным 


отличием от обычного персонального компьютера является отсутствие жесткого диска 


(его заменяет флеш-па\-мять), оптического привода CD/DVD/Blue-ray и встроенного 


блока питания, что позволяет сильно уменьшить размеры процессорного блока. За счет 


умеренного энергопотребления компонентов такие тонкие клиенты не требуют наличия 


мощных и шумных вентиляторов внутри корпуса.


       


       В тонких клиентах осуществляется прошивка <<облегченных>> вариантов 


операционных систем (Microsoft Windows Embedded, HP ThinPro и~т.\,п.) и


десятков различных клиентских программ для подключения к наиболее 


распространенным терминальным серверам, а также к VDI различных производителей. 


Соответственно, цена таких тонких клиентов практически не отличается от цены простых 


офисных компьютеров и, если рассматривать модели известных брендов, находится в 


пределах 300--650~USD. Обычно к данной цене необходимо добавить стоимость 


клиентской лицензии (например, Microsoft Remote Desktop Services CAL) для доступа к 


терминальному серверу или серверу виртуализации, соответственно цена рабочего места 


на базе тонкого клиента увеличится еще на 100--170~USD.


       


       Некоторые универсальные тонкие клиенты, например HP t5745, имеют 


возможность установки в них PCI или PCIe платы, чем можно воспользоваться для 


размещения отечественных про\-грам\-мно-ап\-па\-рат\-ных средств защиты информации, 


предназначенных для персональных компьютеров, и быстро получить реализацию 


технологии тонкого клиента в защищенном исполнении. Такое решение~--- 


       <<QP-тер\-ми\-нал>>~--- предлагается НТП <<Криптософт>>. Очевидно, что 


такой подход наследует все недостатки традиционной кли\-ент-сер\-вер\-ной модели 


вычислений~--- высокую стоимость решения, ориентацию на определенную модель 


импортного тонкого клиента, высокие показатели энергопотребления, сложность 


настройки и обслуживания.


       


       Целесообразным является путь разработки полностью отечественной технологии 


тонкого клиента с изначально встроенными в нее средствами защиты информации. 


Создание любой современной информационной технологии <<с нуля>> представляет 


собой сложнейший процесс с непредсказуемым результатом, поэтому крайне желательно 


наличие реализованных аналогов требуемого решения. В~данном случае полноценным 


прототипом может служить протокол PCoIP и реализующая его технология, предлагаемая 


канадской фирмой Teradici~\cite{2-kor}.





\section{PCoIP как вариант реализации современной технологии тонкого 


клиента}


       


       PCoIP (PC-over-IP, персональный компьютер по протоколу IP)~--- проприетарный 


протокол передачи данных, используемый для передачи компрессированного 


изображения и сигналов различных дополнительных интерфейсов (USB, Audio и~др.)\ 


между сервером виртуализации и терминальным устройством. За счет того, что 


декодирование готового двумерного изображения, сформированного на сервере, является 


достаточно простой задачей, алгоритмы обработки сигнала на стороне клиента могут быть 


эффективно реализованы аппаратным способом на программируемых логических 


интегральных схемах (ПЛИС) типа FPGA (Field-Programmable Gate Array) или на 


недорогих микропроцессорах. Этот подход используется в так называемых аппаратных 


супертонких или нулевых клиентах (hardware zero сlient). Рассматриваемый протокол 


PCoIP преду\-смат\-ри\-ва\-ет сжатие информации об экранном буфере и обеспечивает передачу 


PCoIP-устройствам только данных об изменившихся пикселах. При этом поддерживается 


передача изображения разрешением до $2560\times1600$ и обеспечивается совместимость 


с интерфейсами USB. Требования к пропускной способности сети при использовании 


PCoIP варьируются от 1~Мбит/с при работе с офисными документами до 300~Мбит/с при 


работе с трехмерной графикой в высоком разре-\linebreak\vspace*{-12pt}





\pagebreak





\





\vspace*{-12pt}





\begin{floatingfigure}{67mm} %fig1


       \vspace*{-6pt}


\begin{center}


\mbox{%


\epsfxsize=62mm


\epsfbox{kor-1.eps}


}


\end{center}


\vspace*{-12pt}


\Caption{Teradici APEX 2800 server offload card}


\vspace*{4pt}


\end{floatingfigure}








\noindent


шении или просмотре видео. Самым 


важным свойством для применения данной технологии в защищенных информационных 


системах является изначально заложенная возможность аппаратного шифрования данных 


на PCoIP устройствах.


       


       При использовании технологии PCoIP на стороне центра обработки данных 


(ЦОД) возможно размещение:


       \begin{itemize}


       \item серверных аппаратных решений с использованием специализированного 


ускорителя~--- Teradici APEX 2800 server offload card (рис.~1), обеспечивающего 


аппаратное сжатие на сервере, шифрование и передачу одновременно до 100 сеансов 


изображения удаленного рабочего стола;


       \item программных решений, встроенных в среду виртуализации VMware View 


VDI;


       \item программных решений (Teradici Arch Release~1.0 для Microsoft Remote 


Desktop Services), заменяющих штатный протокол Microsoft RDP (remote desktop protocol)


на протокол PCoIP;


       \item 


        аппаратных решений для отдельных высокопроизводительных рабочих станций, 


установленных в ЦОД, ~--- TERA 2240 host и TERA 2220 host.


       \end{itemize}


         \begin{figure}[b] %fig2


       \vspace*{1pt}


\begin{center}


\mbox{%


\epsfxsize=62mm


\epsfbox{kor-2.eps}


}


\end{center}


\vspace*{-9pt}


\Caption{PCoIP Zero client на базе процессора TERA1100}


\end{figure}











       На стороне защищенного клиента могут быть использованы:


       \begin{itemize}


       \item  аппаратные супертонкие клиенты в виде отдельного устройства (рис.~2);


       \item встроенный в монитор супертонкий клиент PCoIP на базе процессора 


TERA1100 (например, 22$^{\prime\prime}$ монитор Samsung NC220);


       \item программные клиенты для архитектуры $\times86$.


       \end{itemize}


       








       Особый интерес представляет анализ различных вариантов использования 


тонкого клиента PCoIP с целью оценки применимости данного подхода в защищенных 


информационных системах. Далее перечислены основные проблемы, которые необходимо 


решить при создании аналогичной отечественной технологии.


       \begin{enumerate}[1.]


       \item При выполнении на сервере задач пользователей необходимо исключить 


возможность попадания данных одного пользователя к другому, несанкционированного 


понижения степени секретности данных, а также минимизировать влияние выполнения 


задач одних пользователей на выполнение задач других пользователей. Следовательно, 


необходима высокодоверенная среда виртуализации аппаратных ресурсов 


вычислительной техники и доверенная VDI на базе этой среды.


       \item В защищенных информационных системах идентификация пользователя 


должна происходить как минимум с предъявлением физического носителя. 


Следовательно, терминалы должны быть оснащены специализированными доверенными 


устройствами идентификации пользователя, использующими специальные считыватели и 


криптографические методы. 


       \item Информация, передаваемая между сервером виртуализации и тонким 


клиентом, должна быть зашифрована. При разработке технологий тонкого клиента для 


защищенных информационных систем следует ориентироваться на аппаратный способ 


шифрования, обеспечивающий необходимую производительность и сохранность 


ключевой информации.


       \item Использование существующей оптической кабельной сети. Для 


защищенных информационных систем целесообразно использование оптических линий 


передачи данных до конечного рабочего места. В~то же время практически все 


современные импортные терминалы рассчитаны на <<медный>> Ethernet 10/100/1000 


Base-T и для подключения к оптической сети требуют использования промежуточных 


медиаконверторов <<оп\-ти\-ка--медь>>, имеющих сравнимый с тонким клиентом размер и 


электропитание. Следовательно, при разработке отечественного аппаратного тонкого 


клиента необходимо сразу ориентироваться на поддержку двух интерфейсов подключения 


к локальной сети~--- оптическому и медному (витая пара).


       \end{enumerate}


       


     





       При применении технологий тонкого клиента в защищенных информационных 


системах необходимо учитывать и традиционные проблемные моменты, присущие 


технологии в целом.


       


       Не вся информация может быть эффективно передана на терминал, возможны 


проблемы при передаче объемных данных по каналам связи невысокой 


производительности. Например, поток видеоинформации высокой четкости (HD качества)


от терминального 


сервера к терминалу может потребовать полосы пропускания локальной вычислительной сети (ЛВС)


60~Мбит/с. 


       


       В ЛВС не должно быть значительных задержек по времени отклика, в том числе 


из-за использования средств шифрования. В~нормальном режиме работы пользователя 


время отклика должно укладываться в 10--20~мс. 


       


       Возможны проблемы при использовании периферийных USB-устройств 


(принтеры, сканеры), требующих передачи больших объемов информации. 


Распространенный USB~2.0 интерфейс обеспечивает передачу информации со скоростями 


до 480~Мбит/с, а USB~3.0~--- до 5~Гбит/с. И~хотя USB-по\-ток большинством 


терминалов может инкапсулироваться в протоколы сетевого обмена информацией между 


тонким клиентом и сервером виртуализации, ЛВС и адаптер сервера могут не справляться 


с возросшей в сотни раз (по сравнению с обычной архитектурой <<кли\-ент--сер\-вер>>) 


нагрузкой.


       


       Выход из строя (недоступность) сервера виртуализации или других элементов 


VDI-ин\-фра\-струк\-ту\-ры, а также телекоммуникационного оборудования приводит к 


полной невозможности работы пользователя.





\section{Заключение}





       Качественная проработка всех вопросов, начиная от выбора специализированной 


аппаратной платформы сжатия и распаковки изображения, использования шифрования 


при передаче данных и организации взаимодействия с су\-ще\-ст\-ву\-ющей платформой VDI, 


была сделана при создании технологии PCoIP. Используя этот подход, но ориентируясь на 


приемлемую элементную базу, отечественные алгоритмы шифрования, опыт в разработке 


ап\-па\-рат\-но-про\-грам\-мных средств защиты информации, интегрированных с имеющимися 


системными платформами, а также отечественные наработки по виртуализации, вполне 


реально создать оптимизированную под российские условия отечественную технологию 


аппаратного супертонкого клиента и необходимую виртуальную инфраструктуру для его 


работы.





 {\small\frenchspacing


{%\baselineskip=10.8pt


\addcontentsline{toc}{section}{Литература}


\begin{thebibliography}{9}


   





\bibitem{1-kor}


\Au{Соколов И.\,А., Сергеев И.\,В.} Перспективы использования компьютерной 


технологии <<тонкого клиента>> в 


ин\-фор\-ма\-ци\-он\-но-те\-ле\-ком\-му\-ни\-ка\-ци\-он\-ных системах интегрированного 


типа~/\!\!/ Системы и средства информатики.~--- М.: Наука, 2001. Вып.~11. С.~24--37.





 \label{end\stat}


 


\bibitem{2-kor}


PCoIP technology. {\sf http://www.teradici.com/pcoip-technology.php}.








 \end{thebibliography}


}


}


